\section{Results}\label{sec:results}

The results organize around three findings. First, the simulated set-aside price effect is mostly within-auction: the within-auction share is \bneShareIntensNp~percent in non-pharma and \bneShareIntensPh~percent in pharma, and stays in the \withinShareLow--\withinShareHigh~percent range across cost-distribution estimators.\footnote{The DiD-vs-structural relationship is reconciled in Section~\ref{sec:bridge} below.} Second, endogenous SME entry is a partial offset, not the source of the markup: the induced participation response attenuates the within-auction shock, and without it the realized price and welfare cost would be \latentShockRatioNp~percent larger in non-pharmaceuticals and \latentShockRatioPh~percent larger in pharmaceuticals. Third, the policy comparison between bidder exclusion and a \optPrefThreshold{} price preference is conditional on goods characteristics: in thick standardized non-pharma markets the preference rule welfare-dominates robustly; in thin heterogeneous pharma the ranking depends on how the post-policy entrant mix is modeled.

\subsection{Entry response}

Table~\ref{tab:v3_entry_rates} reports the panel. Pre-policy Preg\~ao has \bneNsmePreNp{} SMEs and \bneNnsPreNp{} non-SMEs per auction in non-pharma, \bneNsmePrePh{} and \bneNnsPrePh{} in pharma. Post-policy the counts rotate: \bneNsmePostNp{} SMEs and \nonSmePostNp{} non-SMEs in non-pharma, \bneNsmePostPh{} and \nonSmePostPh{} in pharma. SMEs roughly double in both classes. Non-SMEs fall sharply in both classes. The set-aside therefore changes both the admissible pool and the equilibrium participation response.

\subsection{Structural decomposition of the set-aside}

Table~\ref{tab:v3_bne_decomp} gives the point estimates under the all-bidders regime and endogenous entry, which is the main specification. In non-pharma, the simulated open-regime price is $p_{S_1} = \bneMeanSoneNp$ of the reference price. Holding the SME pool at its pre-period size and restricting the auction to SMEs raises the simulated price to $p_{S_2} = \bneMeanStwoNp$, a within-auction shift of $+\bneEffectIntensNp$. Letting SMEs enter at their post-policy rate pulls the simulated price down to $p_{S_3} = \bneMeanSthreeNp$, a participation-margin offset of $\bneEffectEntryNp$. The net effect is $p_{S_3} - p_{S_1} = +\bneEffectTotalNp$.

The pharmaceutical class follows the same decomposition with larger magnitudes. $p_{S_1} = \bneMeanSonePh$, $p_{S_2} = \bneMeanStwoPh$, $p_{S_3} = \bneMeanSthreePh$, and the total effect is $+\bneEffectTotalPh$. The within-auction shift is $+\bneEffectIntensPh$; the participation offset is $\bneEffectEntryPh$. The larger magnitudes line up with the thinner baseline SME pool in pharmaceuticals (\bneNsmePrePh{} vs \bneNsmePreNp{} SMEs per auction Pre-period, Table~\ref{tab:v3_entry_rates}) and the heavier auction-level heterogeneity (ICC \iccDecompContrastPh{} vs \iccDecompContrastNp{} in Preg\~ao non-SME cells, Table~\ref{tab:v3_uh_variance}).

The within-auction share, defined as $|p_{S_2} - p_{S_1}| / (|p_{S_2} - p_{S_1}| + |p_{S_3} - p_{S_2}|)$, is \bneShareIntensNp~percent in non-pharmaceuticals and \bneShareIntensPh~percent in pharmaceuticals in the main specification. Table~\ref{tab:v3_sensitivity_fc} shows \turnbullShareNp{} and \turnbullSharePh{} percent under the Turnbull NPMLE. Table~\ref{tab:v3_strict_invariance} shows \strictShareNp{} and \strictSharePh{} percent under the strict-primitive-invariance benchmark. Across classes and specifications, the within-auction share sits in the \withinShareLow--\withinShareHigh~percent range. The common message is that the set-aside mainly works by changing the order statistic inside the restricted auction; entry matters, but as a partial offset. Figure~\ref{fig:v3_decomposition} makes this decomposition visible.

\begin{figure}[!htbp]
\centering
\includegraphics[width=0.98\textwidth]{../output/figures/fig_v3_decomposition.pdf}
\caption[Bayes-Nash counterfactual price decomposition.]{\textit{Bayes-Nash counterfactual price decomposition by class, under observed equilibrium entry.} Bars report the simulated second-order statistic $\bar{c}_{(2)}$ normalized by the item reference price, under three counterfactual scenarios. $S_1$ is the open regime at the Pre-period pool (baseline). $S_2$ imposes SME-only at the Pre-period SME pool size (within-auction component isolated). $S_3$ is the SME-only counterfactual at the observed Post-period SME pool (full policy effect with endogenous entry). Annotations report the within-auction component ($S_2-S_1$), the participation-margin component ($S_3-S_2$), and the simulated total effect ($S_3-S_1$). The dotted horizontal line marks the reference price ($p^{\mathrm{ref}} = 1$). Under the main specification, the within-auction component dominates the simulated total effect in both classes; the pharmaceutical decomposition is more model-sensitive (Section~\ref{sec:robustness}). Numbers from Table~\ref{tab:v3_bne_decomp}.}
\label{fig:v3_decomposition}
\end{figure}

\subsection{Endogenous entry as a partial offset}

Table~\ref{tab:v3_apv} makes the role of entry visible. V2 imposes the full set-aside while holding the SME pool fixed at its pre-policy size. Relative to the realized full-set-aside counterfactual V0, V2 is \latentShockTimesNp~percent of V0 in non-pharmaceuticals and \latentShockTimesPh~percent in pharmaceuticals. The fixed-pool simulated effect therefore exceeds the realized one by roughly \latentShockRatioNp~percent in non-pharmaceuticals and \latentShockRatioPh~percent in pharmaceuticals. Both terms in the welfare formula of Section~\ref{sec:welfare} load on this simulated price margin, so the fixed-pool welfare arithmetic would materially overstate the burden of the policy under the model's main specification.

The substantive content is that the within-auction component of the simulated total effect dominates the participation-margin component in non-pharmaceutical markets, and that the participation response attenuates the resulting burden ex post. Reading the participation response as a partial offset to a larger latent within-auction shock is the substantive claim; reading it as a free-standing fiscal-insurance result would overstate the model. The 50--60 percent attenuation is conditional on the policy instrument: under V3 (10 percent price preference), the analogous fixed-pool-vs-realized comparison delivers an attenuation close to zero, since non-SMEs remain in the auction and the participation margin barely moves. The entry response is therefore a property of the full-set-aside instrument, not a free-standing feature of the SME population. Figure~\ref{fig:v3_entry_insurance} shows the size of the simulated offset.

\begin{figure}[!htbp]
\centering
\includegraphics[width=0.98\textwidth]{../output/figures/fig_v3_entry_insurance.pdf}
\caption[Simulated participation-margin offset.]{\textit{Simulated price effect under the realized vs fixed-pool counterfactuals.} Bars report the simulated price effect relative to the open regime, $\Delta p / p^{\mathrm{ref}}$, under two scenarios. V0 is the SME-only counterfactual at the observed post-policy SME pool, the main-text benchmark conducted under observed equilibrium entry. V2 imposes the same SME-only rule at the pre-policy SME pool size. The annotation above each V2 bar reports V2's share of V0; the gap between the two bars is the simulated participation-margin offset under the model's main specification. Numbers from Table~\ref{tab:v3_apv}.}
\label{fig:v3_entry_insurance}
\end{figure}

\subsection{Sources of the departure from earlier estimates}\label{sec:bridge}

Table~\ref{tab:v3_decomp_grid} reports the $2 \times 2$ grid in which each axis is a methodological choice: cost distributions raw versus UH-clean, and entry fixed at the pre-period pool versus endogenous. The clean-and-endogenous row is the main spec; the raw-and-fixed-pool row is the reduced-form-style proxy. The relative gap between them is about \decompGapNpPct~percent in non-pharmaceuticals and about \decompGapPhPct~percent in pharmaceuticals. UH correction accounts for more than half of the gap by changing the right tail of the recovered cost distribution; endogenous entry accounts for the rest by recognizing the policy-induced thickening of the SME pool.

Table~\ref{tab:did_structural_bridge} pushes the same logic further by bridging the DiD coefficient itself to the structural counterfactual. Four mechanical sources of divergence---sample restriction ((a) Preg\~ao only, $c_\epsilon \le \filterCepsUpper$, $n \ge \filterMinBidders$), UH cleaning ((b) Krasnokutskaya scale shrinkage), functional form ((c) second-order statistic vs.\ linear log-mean), and conditioning set ((d) counterfactual vs.\ realized entry)---account for the structural object via additive contributions of $\bridgeSampleNp/\bridgeSamplePh$, $\bridgeUhNp/\bridgeUhPh$, $\bridgeFuncNp/\bridgeFuncPh$, and $\bridgeCondNp/\bridgeCondPh$ in non-pharma and pharma respectively, summing to $\bridgeSumNp/\bridgeSumPh$ on the $p^{\mathrm{ref}}$ normalization. The structural $p_{S_3} - p_{S_1}$ delivered by the BNE simulation is $+\bneEffectTotalNp/+\bneEffectTotalPh$, leaving residuals of $\bridgeResidNp/\bridgeResidPh$ that reflect non-linear interactions between the four sources.

In magnitudes: sample restriction and UH cleaning each carry roughly a quarter to a third of the bridge in non-pharma; functional form (second-order statistic vs. linear log-mean) is the dominant single source in non-pharma at $\bridgeFuncNp$; conditioning set is the smallest contributor in both classes. The pharmaceutical bridge tilts more heavily on UH cleaning ($\bridgeUhPh$, the single largest source) and sample restriction ($\bridgeSamplePh$), reflecting the heavier auction-level heterogeneity documented in Section~\ref{sec:identification}. The 3--4$\times$ gap between the structural and the DiD is therefore predictable in direction and magnitude under the maintained structural assumptions; a misspecified model would not be guaranteed to deliver a gap of either consistent sign or consistent magnitude across two classes. The four-source attribution is interpretive rather than identifying---formal identification would require independent variation in each component (e.g., random sample restrictions, exogenous UH-vs-raw treatment), which the design does not deliver. The table is read as a mechanical bridge, not as a free-standing decomposition theorem.

\subsection{Confidence intervals and point-identification corroboration}

A cluster bootstrap at the auction level with $\bootReps$ replicates (Table~\ref{tab:v3_bootstrap_ci}) gives 95 percent CIs on $p_{S_3} - p_{S_1}$ under the all-bidders regime of $[\bootCILoNp, \bootCIHiNp]$ in non-pharma (bootstrap mean $\bootMeanNp$) and $[\bootCILoPh, \bootCIHiPh]$ in pharma (bootstrap mean $\bootMeanPh$). The bootstrap means lie slightly below the full-sample point estimates of Table~\ref{tab:v3_bne_decomp} ($+\bneEffectTotalNp$ and $+\bneEffectTotalPh$). The gap reflects the interaction between cluster resampling and the nonlinearity of the second-order-statistic functional: with $\bootReps$ cluster-bootstrap replicates, the finite-sample-bias-corrected mean falls below the point estimate; increasing $B$ shrinks the gap and confirms the discrepancy as a finite-$B$ Monte Carlo artifact rather than a structural bias in the point estimate. Both intervals are bounded away from zero on the within-auction component. The bootstrap within-auction-share CI is $[\bootShareCILoNp\%, \bootShareCIHiNp\%]$ in non-pharmaceuticals and $[\bootShareCILoPh\%, \bootShareCIHiPh\%]$ in pharmaceuticals; the Turnbull NPMLE keeps the same qualitative split, with the pharmaceutical within-auction share if anything higher.

\subsection{Entry cost calibration}

Table~\ref{tab:v3_entry_cost} reports the zero-profit implied entry cost per bidder. In non-pharmaceuticals, $\kappa^{\text{SME}} =$ R\$0.55 and $\kappa^{\neg} =$ \entryCostMax{} ($\kappa^{\neg}/\kappa^{\text{SME}} = \entryCostRatioNp$). In pharmaceuticals, $\kappa^{\text{SME}} =$ \entryCostMin{} and $\kappa^{\neg} =$ R\$0.51 ($\kappa^{\neg}/\kappa^{\text{SME}} = \entryCostRatioPh$). The roughly \entryCostRatioNp-to-1 ratio between non-SME and SME entry costs is consistent across classes and plausible given the larger documentation and compliance burden faced by bigger suppliers. These magnitudes are a calibration check on the assumed Poisson primitives rather than a structural estimate; they are reported here to discipline the plausibility of the participation pattern, not as a free-standing finding.\footnote{The $\kappa^k$ values are per-auction zero-profit margins at the entry margin, not fixed costs of preparing a BEC documentation packet (which is amortized over many auctions and is not directly recovered by this calibration). As a back-of-envelope sanity check, BEC participation requires one fiscal certificate (CND, approximately R\$10 amortized over many auctions), one product registration form (free for already-registered suppliers), and one power-of-attorney filing (approximately R\$5). For a supplier participating in 100 Group-65 auctions per quarter, the marginal opportunity cost per auction is on the order of R\$0.10--R\$0.20, consistent with the calibrated $\kappa^{\text{SME}}=$ R\$0.11 in pharma. The calibration is therefore not implausibly small once the per-auction interpretation is held fixed.}

\subsection{Policy comparison and goods characteristics}\label{sec:pareto}

The policy comparison follows directly from the structural decomposition: V3 leaves non-SMEs in the auction, V0 removes them entirely, and the within-auction component identified above prices the difference between the two instruments. Table~\ref{tab:v3_apv} reports a \optPrefThreshold{} SME price preference, under which SME bids are scored with a \optPrefThreshold{} discount for winner selection but the government pays the actual bid. In non-pharma, the preference shifts price by $\optPrefShiftNp$ relative to the open regime; in pharma it shifts price by $\vThreeShiftPh$. The preference-grid simulation (Table~\ref{tab:v3_preference_grid}, $\prefGridB$ draws) gives $\vThreePrefGrid$ in both classes, within Monte Carlo sampling noise of the APV estimates. In both cases the price effect is essentially zero because non-SMEs remain in the auction and continue to discipline the second-order statistic even when an SME wins.

What differs across classes is the stability of the ranking against the full set-aside. In non-pharma, the preference rule dominates under the main equilibrium-selection specification and continues to dominate under the strict-invariance benchmark in Table~\ref{tab:v3_strict_invariance}. In pharma, the preference rule dominates under the main specification but the ranking reverses under strict invariance, where the full set-aside's total price effect rises to $+\siDeltaTotalPh$ and the implied welfare weight needed to justify it falls to \welfWeightSiPh. I interpret that bifurcation as diagnostic of market structure. In thicker, more standardized markets, bidder exclusion is robustly dominated by a moderate preference rule. In thinner, more heterogeneous markets, the ranking is more sensitive because the identity of the induced entrants matters for equilibrium price formation.
